\section{Construction side: redundant archives}\label{sec:limitations}

The necessary conditions above do not imply that two tapes contain independent
information. The following paid normalization rules out that interpretation,
even when the selected words are future-readable and far from the heads.

\begin{theorem}[Redundant-store normalization]\label{thm:redundant}
\notformalized
For every fixed native machine $M$ with $k\ge1$ work tapes there is a
blank-started native machine $\widehat M$ with the same tape count and fixed
enlarged alphabet/control that has the same output on every halting nonempty
input and halts there within $8T_M(X)+20|X|+40$ steps. It halts iff $M$ does.
Fix an input length $P$, any probability law on its halting inputs, and any
absolute configuration time, padding after halt. For any two tapes, their
whole words mutually determine one another. The same is true of their
future-readable words, and of their future-readable words with a common
closed radius-$H$ head window blanked, whenever $2H\le P$.
Head positions are not included in these word variables.
\end{theorem}
\begin{informalproof}
Use fixed product tracks for the original payload, tagged raw archives,
an origin marker, an end marker and a visited flag. Tagged zero differs from
blank. All heads prepare synchronously. Copy the raw input into tagged cells
$[0,P-1]$ on every tape, marking the origin at zero. At the right raw marker
move raw left and work heads right, then continue while raw scans bits. A
turn moves raw from $-1$ to zero while work heads stay at $2P+1$. Copy the
input again into $[2P+1,3P]$. Mark the end at $3P+1$, then return all work
heads to the origin and independently return raw to zero. Counting the bit,
turn and phase transitions gives $6P+4$ preparation steps. There is no length
counter, extra metadata head or preloaded tape. Payload is still blank.

Simulate each original transition on its complete old payload tuple in one
step, preserving the archival tracks. Maintain one optional pending output
bit: on an original append, emit the old pending bit if present and replace
it with the new bit. Otherwise emit nothing. The emitted word followed by
the pending bit equals the original output. This uses three pending states
and never reads output. Mark every old scanned cell as visited, including
each preparation/core step and the stationary wrapper transition detecting
the original halt. Seen cells on each tape form one interval $[l_h,r_h]$
containing $[0,3P+1]$, and every nonblank payload cell is inside it.

At halt, independently on each head, sweep left through visited cells while
blanking payload but retaining the visited flag. At the first unvisited
guard turn right. Sweep right through the whole visited interval, blanking
payload and clearing its visited flag behind the head. At the right unvisited
guard turn left and return to the unique origin. Park completed heads. These
are finite local modes; all heads work in parallel. Archives and markers are
preserved. The seen interval has left endpoint at least $-T_M(X)$ and right
endpoint at most $\max(3P+1,T_M(X))$. A safe joint cleanup cost is
$5T_M(X)+6P+10$. At its end every tape has exactly the same word: the two
archives, public markers and blanks elsewhere. A final synchronized scan
from origin through end successfully reads both archives, then appends the
pending bit and halts. Together these costs are below the stated bound.

Before the public end of preparation, all tape words and heads are identical.
Every then-nonblank cell is an archive or public marker visited by the final
successful scan, so the future-readable words are identical as well.
During the unsynchronized core and cleanup, both tagged archives remain in
each future-readable word because of the final scan. A window of diameter
$2H\le P$ cannot meet both archives, whose nearest cells are distance $P+2$
apart. Thus each whole, future-readable or remote word contains an intact
copy of $X$. During the synchronized tail and after halt, all same-type words
are equal across tapes, since both contents and remaining trajectories agree.

This is also valid for a law mixing inputs in different phases at one absolute
time. A same-type word containing an intact tagged copy determines $X$, hence
the deterministic run and the desired other-tape word at that time. If it
contains no intact copy, it cannot be in the unsynchronized core/cleanup, so
return that same word: the tapes are then synchronized. The public preparation
case is equality as well. These are mathematical decoders, not uncharged
native access routines. Mutual deterministic dependence follows for every
common input law. Divergent output prefixes need not be identical because a
final bit can remain pending; only halting partial-function preservation and
equivalent halting are asserted.
\end{informalproof}

On every timely transpose-correct input the source has at least $P$ successful
steps. During the common first $P$ simulated steps each archived word is an
injective function of the input, yet its conditional information given any
other archived word is zero. The output is still incomplete because of the
pending-bit delay. Ignored future rereads can therefore satisfy large live
marginal-information requirements. Excluding every such normalized transposer
would already settle the fixed-tape-count existence question itself.

An affirmative construction must deliver the required output order through
actual local reads while meeting the return, cut and storage conditions.
Neither large joint support nor a globally available mathematical decoder
pays those accesses. Conversely, an impossibility proof needs a physical
supplier of a small cut, a cheap joint record or another output-sensitive
deficit. None is supplied for unrestricted coupled mutable tapes here.
